\originalchapter{习题课四: 电容与储能}
题源为\href{https://ocw.mit.edu/courses/8-02-physics-ii-electricity-and-magnetism-spring-2007/resources/solving04/}{Problem Solving 4}, PDF第1--7页. 主情境Q1--Q5及球电容考试练习题全部保留. 所有解答为AI独立解答.

\begin{exercise}
同轴导体圆柱半径$a<b$, 长$\ell$, 内柱带$-Q$, 外柱带$+Q$, 忽略端效应.
Q1求$a<r<b$场的大小和方向.
Q2按$\Delta V=V(b)-V(a)$求势差.
Q3求电容.
Q4从$u_E=\eps E^2/2$对全空间积分求能量, 再用$Q,C$表示.
Q5从不带电开始人工把电荷从外柱搬到内柱: (a)已有电荷量$q$时求势差; (b)再搬$\dd q$所做功; (c)积分求搬完$Q$的总功; (d)与Q4能量比较并解释.
\end{exercise}
\begin{solution}
Q1. 高斯柱内电荷负, 所以$\ee=-Q\uv r/(2\pi\eps\ell r)$.
Q2. $\Delta V=-\int_a^b E_r\dd r=Q\ln(b/a)/(2\pi\eps\ell)>0$.
Q3. $C=Q/\Delta V=2\pi\eps\ell/\ln(b/a)$.
Q4. 场在内导体内部和外部为零, 仅间隙贡献:
\[
U=\int_a^b\frac\eps2\frac{Q^2}{(2\pi\eps\ell r)^2}2\pi r\ell\dd r
=\frac{Q^2\ln(b/a)}{4\pi\eps\ell}=\frac{Q^2}{2C}.
\]
Q5. 原件前情境内柱为负, 而本问把“电荷$q$”按正搬运量叙述. 为避免符号混淆, 令$q$为已分离电荷的大小, 两柱势差大小$q/C$. (a)$|\Delta V|=q/C$; 若仍指定$V(b)-V(a)$且内柱为负则$+q/C$. 若真的搬正电荷至内柱, 则该指定势差是$-q/C$. (b)准静态克服电场的功$\dd W=(q/C)\dd q$, 两种一致的符号约定都得到正值. (c)$W=\int_0^Q(q/C)\dd q=Q^2/(2C)$. (d)等于场能, 因为过程没有焦耳耗散和动能变化, 外功全转为场能. 忽略符号约定而同时把负内柱和正内柱混在一个表达式里会导致错误负功.
\end{solution}

\begin{exercise}
考试练习题. 半径$a$实心导体球带$Q$, 外有内半径$b>a$的薄导体球壳带$-Q$. (a)用高斯定律求间隙场并画所用面; (b)求$V(a)-V(b)$; (c)求电容; (d)用体积元$4\pi r^2\dd r$对场能密度积分; (e)核与$C\Delta V^2/2$或$Q^2/(2C)$一致.
\end{exercise}
\begin{solution}
(a)选同心半径$r$高斯球, 得$\ee=Q\uv r/(4\pi\eps r^2)$. 下图为过球心截面, 虚线是$a<r<b$高斯球的截线, 非新的实体导体.
\begin{center}\begin{tikzpicture}[>=Stealth,scale=.8]
\draw[thick](0,0)circle(.7);\draw[thick](0,0)circle(1.7);\draw[dashed,structurecolor](0,0)circle(1.2);
\draw[->](0,0)--(.7,0)node[midway,below]{$a$};\draw[->](0,0)--(0,1.2)node[midway,left]{$r$};\draw[->](0,0)--(-1.2,-1.2)node[midway,left]{$b$};
\node at(.3,.25){$+Q$};\node[right]at(1.7,0){$-Q$};\node at(0,-2.1){同心高斯球};
\end{tikzpicture}\end{center}
(b)$V(a)-V(b)=\int_a^bE_r\dd r=Q(1/a-1/b)/(4\pi\eps)$.
(c)$C=4\pi\eps ab/(b-a)$.
(d)$U=\int_a^b(\eps/2)[Q/(4\pi\eps r^2)]^2 4\pi r^2\dd r=Q^2(1/a-1/b)/(8\pi\eps)$.
(e)将(c)代入$Q^2/(2C)$即得(d); 用$Q=C\Delta V$又得到$C\Delta V^2/2$. 外壳带$-Q$使壳外净包围电荷为零, 且球对称, 所以没有漏算外部场能.
\end{solution}
